\section{Limits of Interpretation and Ethical Scope}

\subsection{Overview}

Because the present framework is designed to admit increasing levels of formal precision and abstraction,
it is necessary to state its interpretative and ethical limits explicitly.
While formalization enhances conceptual clarity,
it can also invite misreadings that conflict with the relational and non-normative intent of the theory.

This section therefore establishes theoretical and ethical boundaries that are integral to the framework itself.

\subsection{Clarification of Scope}

The life-ray framework is not intended to function as a \textit{descriptive}, \textit{diagnostic}, or \textit{person-centered} evaluative instrument.
In particular, it must not be interpreted or employed as:

\begin{itemize}
    \item a diagnostic framework for psychological \textit{classification},
    \item a \textit{prognostic} model for individual development or performance,
    \item a taxonomy for \textit{categorizing} human subjects,
    \item a tool for \textit{institutional selection}, exclusion, or optimization.
\end{itemize}

Evaluation within the framework applies to relational configurations,
system states, and structural dynamics of encounter and memory formation,
not to \textit{persons} as objects of judgment, classification, or assessment.

\subsection{Rationale}

The \emph{Life Ray Theory} is fundamentally non-deterministic and non-normative.
Its purpose is to provide a relational and topological language for understanding how meaning,
memory, and identity emerge through encounter over time.

Historically, the absence of ethical boundary-setting in early foundational theories
has often enabled unintended and irreversible misuse.
This work therefore integrates ethical scope as an explicit and constitutive element
of its theoretical architecture.

Any attempt to translate the framework into classificatory or predictive instruments
applied to \textit{persons} would contradict its core assumptions
and reproduce precisely the reductionist practices the theory seeks to overcome.

\subsection{Relation to Formalization}

The mathematical and computational structures introduced in subsequent parts of the project
serve to support comparative, state-based, and relational evaluation within the model.
Such evaluation is required for formal modeling, simulation,
and computational implementation.

However, these formal structures do not introduce objectivity
in the sense of ranking, norming, or evaluating individuals.
Accordingly, the formal layer should be read as a structural lens
for analyzing relational dynamics, not as an operational metric applied to \textit{persons}.
